\begin{theorem}[歧义壳对 holomorphic 非零谱泛函的透明性]\label{thm:conclusion-ambiguity-shell-holomorphic-nonzero-functional-blindness}
沿用定理 \ref{thm:conclusion-det-fischer-nonzero-jordan-transparent} 的块分解
\[
A_m^{\mathrm{det}}
=
\begin{pmatrix}
N_m & R_m\\
0 & A_m^{\mathrm{ess}}
\end{pmatrix},
\qquad
N_m^{\,m}=0.
\]
设 \(\phi\) 在包含
\[
\sigma(A_m^{\mathrm{det}})\cup \sigma(A_m^{\mathrm{ess}})
\]
的某个开邻域内全纯。则：
\begin{enumerate}
  \item 若 \(\phi(0)=0\)，则
  \[
  \Tr\phi(A_m^{\mathrm{det}})
  =
  \Tr\phi(A_m^{\mathrm{ess}});
  \]
  \item 若 \(\phi(0)=1\)，则
  \[
  \det\phi(A_m^{\mathrm{det}})
  =
  \det\phi(A_m^{\mathrm{ess}}).
  \]
\end{enumerate}
因此，一切仅锚定非零谱的 holomorphic trace/determinant 型不变量都对歧义壳完全失明。
\end{theorem}
\begin{proof}
由 holomorphic functional calculus 的 Cauchy 积分表示，块上三角结构在全纯函数演算下保持不变：
\[
\phi(A_m^{\mathrm{det}})
=
\begin{pmatrix}
\phi(N_m) & *\\
0 & \phi(A_m^{\mathrm{ess}})
\end{pmatrix}.
\]

若 \(\phi(0)=0\)，则由谱映射定理
\[
\sigma(\phi(N_m))=\phi(\sigma(N_m))=\{0\},
\]
故 \(\phi(N_m)\) 为幂零矩阵，从而
\[
\Tr\phi(N_m)=0.
\]
对上式取迹即得第一条。

若 \(\phi(0)=1\)，则
\[
\sigma(\phi(N_m)-I)=\{0\},
\]
故 \(\phi(N_m)-I\) 幂零，因而 \(\phi(N_m)\) 形如 \(I+\) 幂零，满足
\[
\det\phi(N_m)=1.
\]
再对块上三角矩阵取行列式，得到
\[
\det\phi(A_m^{\mathrm{det}})
=
\det\phi(N_m)\,\det\phi(A_m^{\mathrm{ess}})
=
\det\phi(A_m^{\mathrm{ess}}).
\]
证毕。
\end{proof}

\begin{theorem}[非零 Jordan 链的图状唯一提升]\label{thm:conclusion-ambiguity-shell-nonzero-jordan-chain-graph-lift}
\leanverified{paper\_conclusion\_ambiguity\_shell\_nonzero\_jordan\_chain\_graph\_lift}
沿用定理 \ref{thm:conclusion-det-fischer-nonzero-jordan-transparent} 的记号，并取 \(\lambda\neq 0\)。
若
\[
v_0,\dots,v_{\ell-1}\in \CC^{d_{\mathrm{ess}}}
\]
构成 \(A_m^{\mathrm{ess}}\) 上特征值 \(\lambda\) 的一条 Jordan 链，即
\[
(A_m^{\mathrm{ess}}-\lambda I)v_0=0,
\qquad
(A_m^{\mathrm{ess}}-\lambda I)v_j=v_{j-1}
\quad(1\le j\le \ell-1),
\]
则存在唯一向量列
\[
x_0,\dots,x_{\ell-1}\in \CC^{d_{\mathrm{amb}}}
\]
使得
\[
w_j:=(x_j,v_j)\in \CC^{d_{\mathrm{amb}}+d_{\mathrm{ess}}}
\]
构成 \(A_m^{\mathrm{det}}\) 上同长度、同特征值的一条 Jordan 链。更具体地，\(\{x_j\}\) 由递推
\[
(N_m-\lambda I)x_0=-R_m v_0,
\]
\[
(N_m-\lambda I)x_j=x_{j-1}-R_m v_j
\qquad(1\le j\le \ell-1)
\]
唯一确定。
\end{theorem}
\begin{proof}
由于 \(N_m\) 幂零且 \(\lambda\neq 0\)，矩阵 \(N_m-\lambda I\) 可逆。
若要求
\[
w_0=(x_0,v_0)
\]
满足
\[
(A_m^{\mathrm{det}}-\lambda I)w_0=0,
\]
则下块方程正是
\[
(A_m^{\mathrm{ess}}-\lambda I)v_0=0,
\]
而上块方程恰给出
\[
(N_m-\lambda I)x_0=-R_m v_0.
\]

对 \(1\le j\le \ell-1\)，条件
\[
(A_m^{\mathrm{det}}-\lambda I)w_j=w_{j-1}
\]
在下块上等价于
\[
(A_m^{\mathrm{ess}}-\lambda I)v_j=v_{j-1},
\]
在上块上等价于
\[
(N_m-\lambda I)x_j=x_{j-1}-R_m v_j.
\]
由于每一步左侧矩阵均可逆，\(\{x_j\}\) 逐项唯一可解，从而得到唯一提升链。证毕。
\end{proof}

\begin{theorem}[固定超前位数族的 mod-Gaussian 同余]\label{thm:conclusion-kstep-gauge-modgaussian-congruence}
对任意固定整数 \(k,\ell\ge 1\)，设
\[
M_m^{(k)}(\theta):=\EE e^{\theta G_m^{(k)}},
\qquad
M_m^{(\ell)}(\theta):=\EE e^{\theta G_m^{(\ell)}}.
\]
则在 \(\theta=0\) 的某个邻域内存在局部一致收敛
\[
\log\frac{M_m^{(k)}(\theta)}{M_m^{(\ell)}(\theta)}
\longrightarrow
\log\frac{C_k(\theta)}{C_\ell(\theta)}.
\]
特别地，在中心极限定标
\[
\theta=\frac{t}{\sqrt m}
\]
下，对每个固定 \(t\in\RR\) 都有
\[
\frac{\EE\exp\!\bigl(\frac{t}{\sqrt m}G_m^{(k)}\bigr)}
{\EE\exp\!\bigl(\frac{t}{\sqrt m}G_m^{(\ell)}\bigr)}
\longrightarrow 1.
\]
因此，全部固定超前位数族在 Gaussian 尺度上属于同一个 mod-Gaussian 等价类；其差异仅保留在 \(m^0\) 量级的解析前因子中。
\end{theorem}
\begin{proof}
定理 \ref{thm:conclusion-kstep-gauge-cumulant-universality-mdp} 的证明已经给出：存在 \(\varepsilon>0\) 以及解析函数 \(C_k(\theta),C_\ell(\theta)\)，使得在 \(|\theta|<\varepsilon\) 上局部一致成立
\[
M_m^{(k)}(\theta)=C_k(\theta)\lambda(\theta)^m(1+o(1)),
\qquad
M_m^{(\ell)}(\theta)=C_\ell(\theta)\lambda(\theta)^m(1+o(1)).
\]
相除得
\[
\frac{M_m^{(k)}(\theta)}{M_m^{(\ell)}(\theta)}
=
\frac{C_k(\theta)}{C_\ell(\theta)}(1+o(1))
\]
局部一致成立。又因为
\[
M_m^{(k)}(0)=M_m^{(\ell)}(0)=1,
\]
故必有 \(\lambda(0)=1\) 且
\[
C_k(0)=C_\ell(0)=1.
\]
缩小 \(\varepsilon\) 后可设 \(C_k,C_\ell\) 在该邻域内均不为零，于是可选取同一对数支并得到所述对数收敛。

当 \(\theta=t/\sqrt m\) 时，\(\theta\to 0\)，于是
\[
\frac{C_k(t/\sqrt m)}{C_\ell(t/\sqrt m)}\longrightarrow 1,
\]
再结合上面的比值渐近即得
\[
\frac{\EE\exp\!\bigl(\frac{t}{\sqrt m}G_m^{(k)}\bigr)}
{\EE\exp\!\bigl(\frac{t}{\sqrt m}G_m^{(\ell)}\bigr)}
\longrightarrow 1.
\]
证毕。
\end{proof}

\begin{theorem}[周期振幅的剩余类分裂与正能量律]\label{thm:conclusion-periodic-amplitude-residue-energy-splitting}
在定理 \ref{thm:conclusion-ambiguity-shell-periodic-amplitude-obstruction} 的假设下，设
\[
\Psi(n)=\sum_{r=0}^{p-1} c_r e^{2\pi i rn/p}
\]
为其 \(p\)-周期主振幅的 Fourier 展开。则对每个剩余类 \(a\in\{0,\dots,p-1\}\) 都有精确子序列渐近
\[
u_m(np+a)=\rho^{np+a}\Psi(a)+O((\rho-\eta)^{np})
\qquad(n\to\infty).
\]
若记
\[
\overline\Psi:=\frac1p\sum_{a=0}^{p-1}\Psi(a)=c_0,
\]
则周期振幅能量满足
\[
\frac1p\sum_{a=0}^{p-1}\bigl|\Psi(a)-\overline\Psi\bigr|^2
=
\sum_{r=1}^{p-1}|c_r|^2
>0.
\]
因此，等模主极点簇不仅阻止单一常数主项，还把主项精确裂解为 \(p\) 条算术子序列主项，并强制一个严格正的周期振幅能量。
\end{theorem}
\leanverified{paper\_conclusion\_periodic\_amplitude\_residue\_energy\_splitting}
\begin{proof}
由定理 \ref{thm:conclusion-ambiguity-shell-periodic-amplitude-obstruction}，
\[
u_m(n)=\rho^n\Psi(n)+O((\rho-\eta)^n).
\]
将其限制到子序列 \(n\mapsto np+a\) 即得
\[
u_m(np+a)=\rho^{np+a}\Psi(a)+O((\rho-\eta)^{np+a}),
\]
吸收有限因子后即为所述形式。

另一方面，离散 Fourier 正交关系给出
\[
\frac1p\sum_{a=0}^{p-1}e^{2\pi i(r-s)a/p}=\delta_{rs},
\]
故 Parseval 恒等式成立：
\[
\frac1p\sum_{a=0}^{p-1}|\Psi(a)|^2=\sum_{r=0}^{p-1}|c_r|^2.
\]
减去平均项 \(|c_0|^2\) 得
\[
\frac1p\sum_{a=0}^{p-1}\bigl|\Psi(a)-\overline\Psi\bigr|^2
=
\sum_{r=1}^{p-1}|c_r|^2.
\]
若右端为零，则 \(c_r=0\) 对所有 \(r\ge 1\) 成立，从而 \(\Psi\) 为常数，这与定理 \ref{thm:conclusion-ambiguity-shell-periodic-amplitude-obstruction} 中 \(\Psi\) 非常数矛盾。故该能量严格为正。证毕。
\end{proof}

\begin{theorem}[幂零歧义壳的 \texorpdfstring{$\zeta$}{zeta}--熵等谱形变]\label{thm:conclusion-ambiguity-shell-zeta-entropy-isospectral-deformation}
\leanverified{paper\_conclusion\_ambiguity\_shell\_zeta\_entropy\_isospectral\_deformation}
固定任意矩阵 \(A_{\mathrm{ess}}\)。对任意幂零块 \(N,N'\) 与任意耦合块 \(R,R'\)，定义
\[
A=
\begin{pmatrix}
N & R\\
0 & A_{\mathrm{ess}}
\end{pmatrix},
\qquad
A'=
\begin{pmatrix}
N' & R'\\
0 & A_{\mathrm{ess}}
\end{pmatrix}.
\]
记 \(\nu(N)\)、\(\nu(N')\) 分别为 \(N\)、\(N'\) 的幂零指数。则：
\begin{enumerate}
  \item \(A\) 与 \(A'\) 在每个非零特征值处具有完全相同的 Jordan 块型；
  \item 它们的 Artin--Mazur \(\zeta\)-函数、primitive 周期轨道计数与拓扑熵完全一致；
  \item 二者 resolvent 的非零极点集合与极点阶数完全一致，且均由 \(A_{\mathrm{ess}}\) 单独决定；所有仅由歧义壳引入的修正都压缩为有限阶矩阵多项式
  \[
  (I-zN)^{-1}=\sum_{j=0}^{\nu(N)-1}z^jN^j,
  \qquad
  (I-zN')^{-1}=\sum_{j=0}^{\nu(N')-1}z^j(N')^j.
  \]
\end{enumerate}
因此，幂零歧义壳的改变不会生成新的非零谱动力学不变量；它只能改变附着在共同骨架奇性之上的解析前因子，并在零谱方向留下有限记忆多项式影子。
\end{theorem}
\begin{proof}
第一条是定理 \ref{thm:conclusion-det-fischer-nonzero-jordan-transparent} 的直接应用：只要本质块 \(A_{\mathrm{ess}}\) 固定，全部非零 Jordan 数据便固定。

对第二条，块上三角行列式公式给出
\[
\det(I-zA)=\det(I-zN)\det(I-zA_{\mathrm{ess}}),
\qquad
\det(I-zA')=\det(I-zN')\det(I-zA_{\mathrm{ess}}).
\]
由于 \(N,N'\) 幂零，
\[
\det(I-zN)=\det(I-zN')=1,
\]
故
\[
\det(I-zA)=\det(I-zA_{\mathrm{ess}})=\det(I-zA').
\]
于是相应的 Artin--Mazur \(\zeta\)-函数、primitive 周期轨道计数以及由谱半径给出的拓扑熵全部一致。

对第三条，标准块逆公式给出
\[
(I-zA)^{-1}
=
\begin{pmatrix}
(I-zN)^{-1} & z(I-zN)^{-1}R(I-zA_{\mathrm{ess}})^{-1}\\
0 & (I-zA_{\mathrm{ess}})^{-1}
\end{pmatrix},
\]
\[
(I-zA')^{-1}
=
\begin{pmatrix}
(I-zN')^{-1} & z(I-zN')^{-1}R'(I-zA_{\mathrm{ess}})^{-1}\\
0 & (I-zA_{\mathrm{ess}})^{-1}
\end{pmatrix}.
\]
其中 \((I-zN)^{-1}\) 与 \((I-zN')^{-1}\) 都是有限阶矩阵多项式，因此不会产生任何非零极点；全部非零极点只能来自共同因子 \((I-zA_{\mathrm{ess}})^{-1}\)，其极点位置与极点阶数遂完全一致。证毕。
\end{proof}

\begin{corollary}[同步尾部的剩余类极限律]\label{cor:conclusion-sync-tail-residue-class-limit-law}
\leanverified{paper\_conclusion\_sync\_tail\_residue\_class\_limit\_law}
在定理 \ref{thm:conclusion-ambiguity-shell-zeta-sync-splitting} 的 sofic 假设下，若同步尾部满足
\[
\nu_m(T_m>n)=c_{n\bmod p}(m)e^{-\varepsilon_m n}(1+o(1)),
\]
则对每个 \(a\in\{0,\dots,p-1\}\) 都有
\[
\lim_{n\to\infty}
e^{\varepsilon_m(np+a)}\nu_m(T_m>np+a)
=
c_a(m).
\]
特别地，若 \(p>1\) 且 \(\{c_a(m)\}_{a=0}^{p-1}\) 不全相等，则不存在常数 \(c\) 使
\[
\nu_m(T_m>n)=c\,e^{-\varepsilon_m n}(1+o(1)).
\]
\end{corollary}
\begin{proof}
把定理 \ref{thm:conclusion-ambiguity-shell-zeta-sync-splitting} 的尾部渐近限制到子序列 \(n\mapsto np+a\)，便得到
\[
\nu_m(T_m>np+a)=c_a(m)e^{-\varepsilon_m(np+a)}(1+o(1)).
\]
两边乘以 \(e^{\varepsilon_m(np+a)}\) 并令 \(n\to\infty\)，即得所述极限。

若存在常数 \(c\) 满足
\[
\nu_m(T_m>n)=c\,e^{-\varepsilon_m n}(1+o(1)),
\]
则沿每个剩余类 \(a\) 取极限都应得到同一个 \(c\)。这迫使
\[
c_a(m)=c
\qquad(\forall a),
\]
与 \(\{c_a(m)\}\) 不全相等矛盾。证毕。
\end{proof}
